\documentclass[12pt]{article}
\usepackage{amsmath, amssymb, geometry, enumitem, booktabs}
\geometry{margin=1in}
\setlength{\parskip}{0.6em}
\setlength{\parindent}{0pt}

\begin{document}

\begin{center}
\Large \textbf{ECON 300 -- Labour Economics} \\[0.7em]
\large \textbf{Solutions: Practice Problem Set 5} \\[0.3em]

\end{center}

\vspace{1em}

\begin{enumerate}[label=\textbf{Question \arabic*.}]

\item \textbf{Immigration composition}

Total immigrants:
\[
320{,}000
\]

\begin{enumerate}[label=(\alph*)]
    \item Economic immigrants:
    \[
    \frac{160{,}000}{320{,}000}\times 100 = 50\%
    \]

    \[
    \boxed{50\%}
    \]

    \item Family-class immigrants:
    \[
    \frac{96{,}000}{320{,}000}\times 100 = 30\%
    \]

    \[
    \boxed{30\%}
    \]

    \item Refugees / humanitarian immigrants:
    \[
    \frac{64{,}000}{320{,}000}\times 100 = 20\%
    \]

    \[
    \boxed{20\%}
    \]

    \item Economic immigrants rise by \(20\%\):
    \[
    160{,}000(1.20)=192{,}000
    \]

    New total:
    \[
    192{,}000+96{,}000+64{,}000=352{,}000
    \]

    \[
    \boxed{352{,}000}
    \]

    \item New share of economic immigrants:
    \[
    \frac{192{,}000}{352{,}000}\times 100 \approx 54.55\%
    \]

    \[
    \boxed{54.55\%}
    \]
\end{enumerate}

% ------------------------------------------------------------

\item \textbf{Immigration and equilibrium wage/employment}

Demand:
\[
W=32-0.02N
\]

Supply:
\[
W=8+0.01N
\]

\begin{enumerate}[label=(\alph*)]
    \item Initial equilibrium: set demand = supply
    \[
    32-0.02N = 8+0.01N
    \]

    \[
    24=0.03N
    \]

    \[
    N=800
    \]

    Substitute into either curve:
    \[
    W=8+0.01(800)=16
    \]

    \[
    \boxed{N=800,\quad W=16}
    \]

    \item New supply:
    \[
    W=8+0.01(N-120)
    \]

    Simplify:
    \[
    W=8+0.01N-1.2=6.8+0.01N
    \]

    New equilibrium:
    \[
    32-0.02N = 6.8+0.01N
    \]

    \[
    25.2 = 0.03N
    \]

    \[
    N=840
    \]

    Wage:
    \[
    W=6.8+0.01(840)=15.2
    \]

    \[
    \boxed{N=840,\quad W=15.2}
    \]

    \item Change in wage:
    \[
    \Delta W=15.2-16=-0.8
    \]

    Change in employment:
    \[
    \Delta N=840-800=40
    \]

    \[
    \boxed{\Delta W=-0.8,\quad \Delta N=40}
    \]

    \item Interpretation:

    Immigration increases labour supply, which lowers the equilibrium wage and raises employment. In this example, the wage falls by \$0.80 and employment rises by 40 workers.
\end{enumerate}

% ------------------------------------------------------------

\item \textbf{Immigrants and native-born workers as imperfect substitutes}

Demand for native-born workers:
\[
W=40-0.01N_n+0.005N_i
\]

Supply of native-born workers:
\[
W=10+0.02N_n
\]

\begin{enumerate}[label=(\alph*)]
    \item When \(N_i=200\):

    Demand becomes:
    \[
    W=40-0.01N_n+0.005(200)
    \]

    \[
    W=40-0.01N_n+1=41-0.01N_n
    \]

    Set equal to supply:
    \[
    41-0.01N_n = 10+0.02N_n
    \]

    \[
    31=0.03N_n
    \]

    \[
    N_n=1033.33
    \]

    Wage:
    \[
    W=10+0.02(1033.33)=30.67
    \]

    \[
    \boxed{N_n=1033.33,\quad W=30.67}
    \]

    \item When \(N_i=300\):

    Demand becomes:
    \[
    W=40-0.01N_n+0.005(300)
    \]

    \[
    W=40-0.01N_n+1.5=41.5-0.01N_n
    \]

    Set equal to supply:
    \[
    41.5-0.01N_n = 10+0.02N_n
    \]

    \[
    31.5=0.03N_n
    \]

    \[
    N_n=1050
    \]

    Wage:
    \[
    W=10+0.02(1050)=31
    \]

    \[
    \boxed{N_n=1050,\quad W=31}
    \]

    \item Interpretation:

    Since immigrant labour enters positively in the native-born demand equation, immigrants are complements to native-born workers. As immigration rises, both native-born employment and wages increase.
\end{enumerate}

% ------------------------------------------------------------

\item \textbf{Migration decision and present value}

\begin{enumerate}[label=(\alph*)]
    \item Net annual earnings advantage of Country B:
    \[
    58{,}000 - 52{,}000 - 2{,}000 = 4{,}000
    \]

    \[
    \boxed{\$4{,}000}
    \]

    \item Present value over 35 years:
    \[
    PV = 4000\left(\frac{1-(1.04)^{-35}}{0.04}\right)
    \]

    First compute:
    \[
    (1.04)^{-35}\approx 0.253
    \]

    So:
    \[
    PV=4000\left(\frac{1-0.253}{0.04}\right)
    \]

    \[
    PV=4000\left(\frac{0.747}{0.04}\right)
    \]

    \[
    PV=4000(18.675)=74{,}700
    \]

    \[
    \boxed{PV \approx \$74{,}700}
    \]

    \item Subtract moving cost:
    \[
    74{,}700 - 15{,}000 = 59{,}700
    \]

    Since this is positive, migration is worthwhile.

    \[
    \boxed{\text{Migration is worthwhile}}
    \]

    \item If moving cost rises to \$30{,}000:
    \[
    74{,}700 - 30{,}000 = 44{,}700
    \]

    Still positive, so migration is still worthwhile.

    \[
    \boxed{\text{Yes, migration is still worthwhile}}
    \]
\end{enumerate}

% ------------------------------------------------------------

\item \textbf{Immigrant earnings assimilation}

\begin{enumerate}[label=(\alph*)]
    \item Initial earnings gap:
    \[
    44{,}000 - 32{,}000 = 12{,}000
    \]

    \[
    \boxed{\$12{,}000}
    \]

    \item Earnings gaps:

    After 5 years:
    \[
    47{,}000 - 39{,}000 = 8{,}000
    \]

    After 10 years:
    \[
    50{,}000 - 46{,}000 = 4{,}000
    \]

    After 15 years:
    \[
    54{,}000 - 52{,}000 = 2{,}000
    \]

    \[
    \boxed{8{,}000,\quad 4{,}000,\quad 2{,}000}
    \]

    \item Gap closed after 15 years:
    \[
    12{,}000 - 2{,}000 = 10{,}000
    \]

    \[
    \boxed{\$10{,}000}
    \]

    \item Percentage of initial gap closed:
    \[
    \frac{10{,}000}{12{,}000}\times 100 = 83.33\%
    \]

    \[
    \boxed{83.33\%}
    \]
\end{enumerate}

% ------------------------------------------------------------

\item \textbf{Point system calculation}

\begin{enumerate}[label=(\alph*)]
    \item Total points:
    \[
    21+20+11+8+0+7=67
    \]

    \[
    \boxed{67}
    \]

    \item Since passing score is 67:
    \[
    \boxed{\text{Yes, the applicant qualifies}}
    \]

    \item If language points improve by 5:
    \[
    67+5=72
    \]

    \[
    \boxed{72}
    \]

    \item Since the applicant already has 67, additional points needed:
    \[
    \boxed{0}
    \]
\end{enumerate}


\item \textbf{Labour force, unemployment rate, LFPR, and employment rate}

\begin{itemize}
    \item Population \(=2{,}400{,}000\)
    \item Employed \(=1{,}560{,}000\)
    \item Unemployed \(=120{,}000\)
\end{itemize}

\begin{enumerate}[label=(\alph*)]
    \item Labour force:
    \[
    1{,}560{,}000+120{,}000=1{,}680{,}000
    \]

    \[
    \boxed{1{,}680{,}000}
    \]

    \item Unemployment rate:
    \[
    \frac{120{,}000}{1{,}680{,}000}\times 100 = 7.14\%
    \]

    \[
    \boxed{7.14\%}
    \]

    \item Labour force participation rate:
    \[
    \frac{1{,}680{,}000}{2{,}400{,}000}\times 100 = 70\%
    \]

    \[
    \boxed{70\%}
    \]

    \item Employment rate:
    \[
    \frac{1{,}560{,}000}{2{,}400{,}000}\times 100 = 65\%
    \]

    \[
    \boxed{65\%}
    \]
\end{enumerate}

% ------------------------------------------------------------

\item \textbf{Labour force and unemployment rate}

\begin{itemize}
    \item Labour force \(=1{,}800{,}000\)
    \item Unemployment rate \(=7.5\%\)
\end{itemize}

\begin{enumerate}[label=(\alph*)]
    \item Number unemployed:
    \[
    0.075(1{,}800{,}000)=135{,}000
    \]

    \[
    \boxed{135{,}000}
    \]

    \item Number employed:
    \[
    1{,}800{,}000-135{,}000=1{,}665{,}000
    \]

    \[
    \boxed{1{,}665{,}000}
    \]

    \item Employment rises by 45{,}000 and unemployment falls by 45{,}000:

    New unemployed:
    \[
    135{,}000-45{,}000=90{,}000
    \]

    New unemployment rate:
    \[
    \frac{90{,}000}{1{,}800{,}000}\times 100 = 5\%
    \]

    \[
    \boxed{5\%}
    \]
\end{enumerate}

% ------------------------------------------------------------

\item \textbf{Monthly labour market transitions}

\begin{enumerate}[label=(\alph*)]
    \item Net monthly change in unemployment:
    \[
    180{,}000 - 210{,}000 = -30{,}000
    \]

    \[
    \boxed{-30{,}000}
    \]

    \item New stock of unemployed:
    \[
    960{,}000 - 30{,}000 = 930{,}000
    \]

    \[
    \boxed{930{,}000}
    \]

    \item Unemployment rate:
    \[
    \frac{930{,}000}{16{,}000{,}000}\times 100 = 5.81\%
    \]

    \[
    \boxed{5.81\%}
    \]
\end{enumerate}

% ------------------------------------------------------------

\item \textbf{Incidence and duration}

\[
UR = \text{Incidence} \times \text{Duration}
\]

\begin{enumerate}[label=(\alph*)]
    \item Region A:
    \[
    2.5\%\times 3 = 7.5\%
    \]

    \[
    \boxed{7.5\%}
    \]

    \item Region B:
    \[
    1.5\%\times 5 = 7.5\%
    \]

    \[
    \boxed{7.5\%}
    \]

    \item Both regions have the same unemployment rate:
    \[
    \boxed{7.5\%}
    \]

    \item Interpretation:

    The same unemployment rate can come from many short spells (high incidence, short duration) or fewer long spells (low incidence, long duration).
\end{enumerate}

% ------------------------------------------------------------

\item \textbf{Youth vs adults}

\begin{enumerate}[label=(\alph*)]
    \item Youth:
    \[
    6.0\%\times 2.0 = 12.0\%
    \]

    Adults:
    \[
    2.0\%\times 2.8 = 5.6\%
    \]

    So the approximation holds exactly here.

    \[
    \boxed{\text{Yes, it holds for both groups}}
    \]

    \item Higher incidence:
    \[
    6.0\% > 2.0\%
    \]

    \[
    \boxed{\text{Youth}}
    \]

    \item Longer duration:
    \[
    2.8 > 2.0
    \]

    \[
    \boxed{\text{Adults}}
    \]

    \item Interpretation:

    Youth have high unemployment because they enter and leave jobs frequently. Their unemployment spells are short, but incidence is high.
\end{enumerate}

% ------------------------------------------------------------

\item \textbf{Labour force, employment, and unemployment}

\begin{itemize}
    \item Population \(=5{,}000{,}000\)
    \item LFPR \(=66\%\)
    \item Employment rate \(=61\%\)
\end{itemize}

\begin{enumerate}[label=(\alph*)]
    \item Labour force:
    \[
    0.66(5{,}000{,}000)=3{,}300{,}000
    \]

    \[
    \boxed{3{,}300{,}000}
    \]

    \item Number employed:
    \[
    0.61(5{,}000{,}000)=3{,}050{,}000
    \]

    \[
    \boxed{3{,}050{,}000}
    \]

    \item Number unemployed:
    \[
    3{,}300{,}000-3{,}050{,}000=250{,}000
    \]

    \[
    \boxed{250{,}000}
    \]

    \item Unemployment rate:
    \[
    \frac{250{,}000}{3{,}300{,}000}\times 100 \approx 7.58\%
    \]

    \[
    \boxed{7.58\%}
    \]
\end{enumerate}

% ------------------------------------------------------------

\item \textbf{Discouraged workers and the unemployment rate}

Initial:
\begin{itemize}
    \item Labour force \(=3{,}000{,}000\)
    \item Unemployed \(=180{,}000\)
\end{itemize}

\begin{enumerate}[label=(\alph*)]
    \item Initial unemployment rate:
    \[
    \frac{180{,}000}{3{,}000{,}000}\times 100 = 6\%
    \]

    \[
    \boxed{6\%}
    \]

    \item New unemployed:
    \[
    180{,}000 + 30{,}000 = 210{,}000
    \]

    \[
    \boxed{210{,}000}
    \]

    \item New labour force:
    \[
    3{,}000{,}000 - 40{,}000 + 30{,}000 = 2{,}990{,}000
    \]

    \[
    \boxed{2{,}990{,}000}
    \]

    \item New unemployment rate:
    \[
    \frac{210{,}000}{2{,}990{,}000}\times 100 \approx 7.02\%
    \]

    \[
    \boxed{7.02\%}
    \]

    \item Interpretation:

    The unemployment rate does not capture all hardship because discouraged workers who stop searching are removed from the labour force and no longer counted as unemployed.
\end{enumerate}

% ------------------------------------------------------------

\item \textbf{Long-term unemployment}

\begin{enumerate}[label=(\alph*)]
    \item Country X:
    \[
    6.0\%\times 10\% = 0.6\%
    \]

    \[
    \boxed{0.6\%}
    \]

    \item Country Y:
    \[
    6.0\%\times 35\% = 2.1\%
    \]

    \[
    \boxed{2.1\%}
    \]

    \item More serious long-term unemployment problem:
    \[
    \boxed{\text{Country Y}}
    \]

    \item Interpretation:

    Even with the same overall unemployment rate, Country Y has a much larger share of workers unemployed for long periods, which suggests more serious structural labour market problems.
\end{enumerate}

% ------------------------------------------------------------

\item \textbf{Canada labour market indicators}

\begin{itemize}
    \item Unemployment rate \(=5.8\%\)
    \item LFPR \(=65.5\%\)
    \item Population \(=31{,}000{,}000\)
\end{itemize}

\begin{enumerate}[label=(\alph*)]
    \item Labour force:
    \[
    0.655(31{,}000{,}000)=20{,}305{,}000
    \]

    \[
    \boxed{20{,}305{,}000}
    \]

    \item Number unemployed:
    \[
    0.058(20{,}305{,}000)=1{,}177{,}690
    \]

    \[
    \boxed{1{,}177{,}690}
    \]

    \item Number employed:
    \[
    20{,}305{,}000-1{,}177{,}690=19{,}127{,}310
    \]

    \[
    \boxed{19{,}127{,}310}
    \]

    \item Employment rate:
    \[
    \frac{19{,}127{,}310}{31{,}000{,}000}\times 100 \approx 61.70\%
    \]

    \[
    \boxed{61.70\%}
    \]
\end{enumerate}

\end{enumerate}

\end{document}